% Source: source/普通高考/2012/2012安徽理.pdf, page 1, question 3.
% pdfimages -list reports no embedded bitmap; the program flowchart is vector artwork.
% Grid evidence: tmp/redraw_flowchart_2012_anhui_science/grid/q3_grid.jpg.
% Nodes: 开始; x=1,y=1; x≤4?; x=2x,y=y+1; 输出 y; 结束.
% Directed edges: 开始→x=1,y=1→x≤4?; 是→x=2x,y=y+1→shared loop entry→x≤4?;
% 否→输出 y→结束. The loop return arrow enters the stem above the decision;
% no line is shared with the output branch.
\begin{tikzpicture}[
  >=Stealth,
  line width=0.55pt,
  line cap=round,
  line join=round,
  every node/.style={font=\normalsize,anchor=center,align=center,
    execute at begin node={\setlength{\parindent}{0pt}}},
  flowbox/.style={draw,minimum height=.68cm,inner xsep=4pt,inner ysep=2pt,
    anchor=center,align=center},
  terminal/.style={flowbox,rounded corners=.18cm,minimum width=1.30cm,
    minimum height=.72cm},
  io/.style={flowbox,shape=trapezium,trapezium left angle=72,
    trapezium right angle=108,minimum height=.78cm},
  decision/.style={flowbox,shape=diamond,aspect=2.8,inner sep=2pt},
  flowarrow/.style={->,line width=0.55pt},
]
  \node[terminal] (start) at (0,7.35) {开始};
  \node[flowbox,minimum width=2.35cm] (init) at (0,6.10) {$x=1,y=1$};
  \coordinate (loopjoin) at (0,5.52);
  \node[decision,minimum width=3.35cm,minimum height=.94cm] (test) at (0,4.62)
    {$x\leq4?$};
  \node[flowbox,minimum width=3.55cm] (update) at (3.75,4.62)
    {$x=2x,y=y+1$};
  \node[io,minimum width=1.85cm] (output) at (0,3.15) {输出 $y$};
  \node[terminal] (finish) at (0,1.88) {结束};

  \draw[flowarrow] (start.south) -- (init.north);
  \draw (init.south) -- (loopjoin);
  \draw[flowarrow] (loopjoin) -- (test.north);
  \draw[flowarrow] (test.east) -- (update.west);
  \draw[flowarrow] (test.south) -- node[right=2pt] {否} (output.north);
  \draw[flowarrow] (output.south) -- (finish.north);

  % The affirmative update returns above the decision into the shared stem.
  \draw (update.north) -- (3.75,5.52);
  \draw[flowarrow] (3.75,5.52) -- (loopjoin);
  \node[anchor=south west,inner sep=0pt] at (1.05,4.98) {是};

  \path[use as bounding box] (-2.10,1.45) rectangle (5.50,7.78);
\end{tikzpicture}
